% \subsection{Case $r=1$} 
% In the case of $r=1$, which serves as a natural starting point for introducing our framework, each vertex must belong to at least $k$ $s$-cliques in the $k$-(1,s)-nucleus subgraph. To obtain the nucleus decomposition, we use a peeling approach that iteratively removes vertices with the fewest $s$-cliques until all vertices satisfy the $k$-(1,s)-nucleus condition.

% We observe the following lemma regarding the update of the Succinct Clique Tree (CPI) after removing a vertex:

% \begin{lemma}[CPI update under vertex removal]\label{lem:path_vertex_remove}
% Let $\TPath = {V_{p}(\TPath), V_{h}(\TPath)}$ be a clique path in the CPI of a graph $G$, $V_{p}(\TPath)$ be the set of pivot vertices on $\TPath$, and $V_{h}(\TPath)$ be the set of hold vertices on $\TPath$. After removing a vertex $v$ from $G$, the update to $\TPath$ is as follows.
% \begin{enumerate}
% \item If $v \in V_{h}(\TPath)$, then the entire path $\TPath$ is removed from the CPI.
% \item If $v \in V_{p}(\TPath)$, then $v$ is removed from $\TPath$ while the rest of the path remains unchanged.
% \end{enumerate}
% \end{lemma}

% \begin{proof}
% By Lemma~\ref{lem:clique_count}, every $k$-clique encoded by a path
% $\TPath{}$ must contain all hold vertices on that path, whereas pivot
% vertices are optional. If the removed vertex $v$ lies in the hold set
% $V_h(\TPath{})$, no encoded $k$-clique can survive, so the whole path
% vanishes. If $v\in V_p(\TPath{})$, only those cliques that actually
% include $v$ disappear; the cliques formed by the unchanged hold set and
% the remaining pivots are still uniquely represented by the shortened
% path, so we remove $v$ and keep the rest of $\TPath{}$ intact.
% \end{proof}

% \begin{example}
% Let us consider the CPI shown in Figure~\ref{fig:CPI_example}. For the CPI shown in Figure~\ref{fig:CPI_example}, consider the clique path $\TPath{1}$ with
% $V_{p}(\TPath{1}) = \{v_4, v_5\}$ and $V_{h}(\TPath{1}) = \{v_1, v_6\}$. 
% If we remove a hold vertex $v_1$ from the underlying graph, then according to Lemma~\ref{lem:path_vertex_remove}, the entire path $\TPath{1}$ will be removed from the CPI. On the other hand, if we remove a pivot vertex $v_4$ from the underlying graph, then according to Lemma~\ref{lem:path_vertex_remove}, we will remove $v_4$ from $\TPath{1}$, and the updated path will be $\TPath{1}'$ with $V_{p}(\TPath{1}') = \{v_5\}$ and $V_{h}(\TPath{1}') = \{v_1, v_6\}$.
% \end{example}

% \begin{algorithm}[t]
\small
\small
\caption{\textsc{UpdateIndex}(r=1)}
\label{alg:RemoveNode}
\KwIn{clique path $\TPath{} = \{V_{H},V_{P}\}$, vertices to remove $V_{\mathrm{rm}}$}
\KwOut{subpaths induced after deleting $V_{\mathrm{rm}}$}
\SetKwFunction{RemoveNode}{RemoveNode}
\SetKwFunction{CountPreVertices}{$s$-CliqueCount}
\SetKwProg{Fn}{Function}{:}{}


\lIf{$|V_{rm} \cap V_{H}(\TPath)| != 0$}{
    \Return
}
% output path $(V_{Piv}(\TPath) \setminus V, V_{Hold}(\TPath))$\;
\textbf{output path }$\{ V_{H}(\TPath),V_{P}(\TPath) \setminus V_{\mathrm{rm}}\}$\;
    %

\end{algorithm}


% \paragraph{Algorithm Description.}
% In Algorithm~\ref{alg:peeling_r1}, we first count the number of $s$-cliques for each vertex in the tree $\tree$ using lemma~\ref{lem:clique_count} in line 2.
% Subsequently, we use a min-heap to store the number of $s$-cliques for each vertex, which allows us to efficiently retrieve the vertex with the minimum $s$-clique count in line 4.
% In each iteration, we extract the vertex $v_{min}$ with the minimum $s$-clique count from the min-heap and update the current minimum core value $c_{min}$ in line 6.
% Then, we set the core value of $v_{min}$ to $c_{min}$ and call the function $\RemoveNode$ to remove $v_{min}$ from each path containing it in line 9.
% In the function $\RemoveNode$, we update the CPI based on Lemma~\ref{lem:path_vertex_remove} in lines 12-20.
% Finally, we return the core values of all vertices in line 22.

% \paragraph{Complexity Analysis.}
% \todo{}

% The initialization first computes $ \textsc{s\textrm{-}CliqueCount}(\mathcal{T},1,s)$ on the CPI $\mathcal{T}$ to obtain $\textit{counts}_s(v)$ for every $v$(Line 1). Each $s$-clique contributes a constant amount of work along a clique path in $\mathcal{T}$, hence the time complexity of this step is $O(\alpha^2|CPI|)$ and its space usage is $O(n+\Sigma_s)$. Building the min-heap $H_{\min}$ over the $n$ vertices costs $O(n)$ time(Line 2).

% During the main loop, each iteration performs one \textsc{ExtractMin} on $H_{\min}$(Line 5), which costs $O(\log n)$. The algorithm then removes the current vertex from all clique paths that contain it and updates counts of the affected vertices(Line 10). Each removal touches every $s$-clique that contains the removed vertex exactly once. Over the whole execution, the total number of affected $s$-cliques equals $\Sigma_s$. For each affected $s$-clique, \textsc{RemoveNode}(Line 15-21) modifies at most $s$ items on a path in $\mathcal{T}$ and triggers at most $s$ heap decrease--key operations when updating the counts of the remaining vertices. Therefore, the total time spent in \textsc{RemoveNode} and heap updates is
% $O(\Sigma_s + \Sigma_s \log n)$, which simplifies to $O(\Sigma_s \log n)$.

% Putting the pieces together, the overall time complexity of Algorithm~6 is
% \[
% O\big(\underbrace{\alpha^2|CPI|}_{\textsc{s\textrm{-}CliqueCount}} + \underbrace{n}_{\text{heap build}} + 
% \underbrace{n\log n}_{\text{all extracts}} + \underbrace{\Sigma_s \log n}_{\text{all updates}}\big)
% = O\big((n+\Sigma_s)\log n + \Sigma_s\big).
% \]
% When $s$ is treated as a constant and $|K_s|$ is large, this is $O\big((n+s|K_s|)\log n\big)$. The space complexity is dominated by the CPI and the heap, which yields $O(n+\Sigma_s)$ overall.


% \begin{example}
%     \todo{}
% \end{example}

\subsection{Case $r=1$} 
\label{sec:r1}
% Within the base framework (Sec.~\ref{sec:base_framework}), the case $r{=}1$ specializes the removal primitive to \emph{vertex removal}. A vertex $v$ belongs to the $k$-(1,$s$)-nucleus iff it participates in at least $k$ $s$-cliques. Peeling removes vertices with the current minimum $s$-clique count and updates the CPI and counters without enumerating $s$-cliques. The only $r{=}1$-specific ingredient is how a vertex removal transforms an CPI path, which we formalize next.
We begin with the basic case of $r=1$ (removing vertices in \cpi). We establish the following lemma after removing a vertex:

% \begin{lemma}[Vertex Removal]\label{lem:path_vertex_remove}
% Let $\TPath = \{V_{p}(\TPath), V_{h}(\TPath)\}$ be a clique path in the \cpi\ of a graph $G$, 
% where $V_{p}(\TPath)$ and $V_{h}(\TPath)$ denote the sets of pivot and hold vertices on $\TPath$, respectively. 
% %
% After removing a vertex $v$ from $G$, the corresponding update to $\TPath$ is as follows:
% % If a vertex $v$ is removed from $G$, the corresponding clique path $\TPath$ in the \cpi\ satisfies the following properties:
% \begin{enumerate}
%     \item If $v \in V_{h}(\TPath)$, the entire path $\TPath$ is removed from the \cpi.
%     \item If $v \in V_{p}(\TPath)$, the vertex $v$ is removed from $\TPath$, while the remaining path stays unchanged.
% \end{enumerate}
% \end{lemma}
\todo{}
\begin{theorem}[Vertex Removal]\label{lem:path_vertex_remove}
Let $\TPath$ be a clique path in the \cpi\ of a graph $G$, with hold and pivot
sets denoted by $V_h(\TPath)$ and $V_p(\TPath)$, respectively.
For any removed vertex $v \in \TPath$, the updated $\TPath'$ satisfies:

\begin{enumerate}[leftmargin=*]
    \item if $v \in V_h(\TPath)$, $\TPath' = \emptyset$;
    \item if $v \in V_p(\TPath)$, $\TPath' = \{(V_p(\TPath) \setminus \{v\},\, V_h(\TPath))\}$.
\end{enumerate}


% For all $s$-cliques encoded by $\TPath$ that do not contain $v$, they remain encoded by $\TPath'$.
\end{theorem}


\begin{proof}
By Lemma~\ref{lem:clique_count}, every $k$-clique encoded by a path
$\TPath{}$ must contain all hold vertices on that path, whereas pivot
vertices are optional. If the removed vertex $v$ lies in the hold set
$V_h(\TPath{})$, no encoded $k$-clique can survive, so the whole path
vanishes. If $v\in V_p(\TPath{})$, only those cliques that actually
include $v$ disappear; the cliques formed by the unchanged hold set and
the remaining pivots are still uniquely represented by the shortened
path, so we remove $v$ and keep the rest of $\TPath{}$ intact.
\end{proof}

\begin{example}
Let us consider the \cpi shown in Figure~\ref{fig:pruning_example}. For the \cpi shown in Figure~\ref{fig:pruning_example}, consider the clique path $\TPath{2}$ with
$V_{p}(\TPath{2}) = \{v_4, v_5\}$ and $V_{h}(\TPath{2}) = \{v_1, v_6\}$. 
If we remove a hold vertex $v_1$ from the underlying graph, then according to Lemma~\ref{lem:path_vertex_remove}, the entire path $\TPath{2}$ will be removed from the \cpi. On the other hand, if we remove a pivot vertex $v_4$ from the underlying graph, then according to Lemma~\ref{lem:path_vertex_remove}, we will remove $v_4$ from $\TPath{2}$, and the updated path will be $\TPath{2}'$ with $V_{p}(\TPath{2}') = \{v_5\}$ and $V_{h}(\TPath{2}') = \{v_1, v_6\}$.
\end{example}

\begin{algorithm}[t]
\small
\small
\caption{\textsc{UpdateIndex}(r=1)}
\label{alg:RemoveNode}
\KwIn{clique path $\TPath{} = \{V_{H},V_{P}\}$, vertices to remove $V_{\mathrm{rm}}$}
\KwOut{subpaths induced after deleting $V_{\mathrm{rm}}$}
\SetKwFunction{RemoveNode}{RemoveNode}
\SetKwFunction{CountPreVertices}{$s$-CliqueCount}
\SetKwProg{Fn}{Function}{:}{}


\lIf{$|V_{rm} \cap V_{H}(\TPath)| != 0$}{
    \Return
}
% output path $(V_{Piv}(\TPath) \setminus V, V_{Hold}(\TPath))$\;
\textbf{output path }$\{ V_{H}(\TPath),V_{P}(\TPath) \setminus V_{\mathrm{rm}}\}$\;
    %

\end{algorithm}

% \begin{algorithm}[t]
\small
% \caption{\textsc{RemoveNode} (r=1 path update)}
% \label{alg:RemoveNode}
% \KwIn{CPI $\tree$, clique path $\TPath$, vertex $v$, clique size $s$}
% \KwOut{Core value $\kappa_s(v)$ for every vertex $v$}
% \SetKwFunction{RemoveNode}{RemoveNode}
% \SetKwFunction{CountPreVertices}{$s$-CliqueCount}
% \SetKwProg{Fn}{Function}{:}{}


% \If{$v$ is hold in $\TPath$}{ 
%     remove $\TPath$ from $\tree$\;
% }
% \ElseIf{$v$ is pivot in $\TPath$}{
%     remove $v$ from $\TPath$\;
%     \If{$|\TPath| < s$}{
%         remove $\TPath$ from $\tree$\;
%     }
% }
% % }

\stitle{Batched Updates.}
% 我们可以很容易地将算法扩展到批量删除顶点的情况。给定一个顶点删除集合$V_{\mathrm{rm}}$和一个包含任何被删除顶点的路径$\TPath$，我们可以通过以下方式更新$\TPath$:
In each iteration of the peeling process, many $r$-cliques are removed. 
To avoid redundant computation, we adopt a batched strategy that merges repeated updates. 
For vertex removals, the algorithm can be easily extended to handle batches. 
Given a set of vertices to be removed $V_{\mathrm{rm}}$ and a path $\TPath$ containing any of them, we update $\TPath$ as follows:
\begin{enumerate}[leftmargin=*]
    \item $\TPath' = \emptyset$ if $|V_h(\TPath) \cap V_{\mathrm{rm}}| > 0$;
    \item $\TPath' = \{(V_p(\TPath) \setminus V_{\mathrm{rm}},\, V_h(\TPath))\}$ if $|V_h(\TPath) \cap V_{\mathrm{rm}}| = 0$.
\end{enumerate}



% \end{algorithm}
\stitle{Algorithm.} 
Algorithm~\ref{alg:RemoveNode} implements the path update procedure described in Lemma~\ref{lem:path_vertex_remove},
% 给定一个顶点删除集合$V_{\mathrm{rm}}$和一个包含任何被删除顶点的路径$\TPath$，我们根据上述规则更新$\TPath$。
Given a set of vertices to be removed $V_{\mathrm{rm}}$ and a path $\TPath$ containing any of the removed vertices, we update $\TPath$ according to the rules above.


\stitle{Complexity.}
Let a clique path $\TPath = \{V_p(\TPath), V_h(\TPath)\}$.
For a single path update, the running time of algorithm \ref{alg:RemoveNode} is $O\!\big(|V_h(\TPath)|+|V_p(\TPath)|\big)$
and the extra space is $O(1)$.